\documentclass[11pt]{article}

\usepackage[T1]{fontenc}
\usepackage{amsmath,amssymb,amsthm,mathtools,bm}
\usepackage{newtxtext,newtxmath}
\usepackage[letterpaper,margin=1in,headheight=15pt]{geometry}
\usepackage{microtype}
\usepackage{booktabs,longtable,tabularx,array}
\usepackage[most]{tcolorbox}
\usepackage{xcolor}
\usepackage{enumitem}
\usepackage{fancyhdr}
\usepackage{titlesec}
\usepackage{graphicx}
\usepackage{listings}
\usepackage{xurl}
\usepackage{hyperref}
\usepackage[nameinlink,noabbrev]{cleveref}
\usepackage{bookmark}
\usepackage{needspace}

\definecolor{ink}{HTML}{18212B}
\definecolor{accent}{HTML}{1E5A7A}
\definecolor{accenttwo}{HTML}{7A3E1E}
\definecolor{pale}{HTML}{F3F7F9}
\definecolor{palewarm}{HTML}{FBF6F1}
\definecolor{rulegray}{HTML}{CBD5DB}
\definecolor{codegray}{HTML}{F5F5F5}

\hypersetup{
  colorlinks=true,
  linkcolor=accent,
  citecolor=accenttwo,
  urlcolor=accent,
  pdftitle={Reversing x + W(x): Exact Reduction and Logarithmic Transseries},
  pdfauthor={OpenAI},
  pdfsubject={Lambert W, r-Lambert function, asymptotic inversion, logarithmic transseries},
  pdfkeywords={Lambert W, r-Lambert W, Lagrange inversion, transseries, asymptotic expansion}
}

\pagestyle{fancy}
\fancyhf{}
\fancyhead[L]{\small\textsc{Reversing $x+W(x)$}}
\fancyhead[R]{\small\textsc{Logarithmic transseries}}
\fancyfoot[C]{\thepage}
\renewcommand{\headrulewidth}{0.35pt}

\titleformat{\section}
  {\Large\bfseries\color{ink}}{\thesection}{0.65em}{}
\titleformat{\subsection}
  {\large\bfseries\color{accent}}{\thesubsection}{0.65em}{}
\titleformat{\subsubsection}
  {\normalsize\bfseries\color{accenttwo}}{\thesubsubsection}{0.65em}{}
\titlespacing*{\section}{0pt}{2.4ex plus .5ex minus .2ex}{1.2ex}
\titlespacing*{\subsection}{0pt}{2.0ex plus .4ex minus .2ex}{0.8ex}

\setlist{itemsep=0.25em,topsep=0.5em}
\setlength{\parindent}{0pt}
\setlength{\parskip}{0.52em}
\allowdisplaybreaks[3]

\newtheoremstyle{clean}
  {8pt}{8pt}{\itshape}{}% space, body font, indent
  {\bfseries\color{ink}}{.}{0.55em}{}
\theoremstyle{clean}
\newtheorem{theorem}{Theorem}[section]
\newtheorem{proposition}[theorem]{Proposition}
\newtheorem{lemma}[theorem]{Lemma}
\newtheorem{corollary}[theorem]{Corollary}
\newtheorem{definition}[theorem]{Definition}
\newtheorem{remark}[theorem]{Remark}
\newtheorem{example}[theorem]{Example}

\newtcolorbox{mainbox}[1][]{
  enhanced,
  breakable,
  colback=pale,
  colframe=accent,
  boxrule=0.7pt,
  arc=2pt,
  left=9pt,right=9pt,top=8pt,bottom=8pt,
  title={#1},
  fonttitle=\bfseries
}
\newtcolorbox{methodbox}[1][]{
  enhanced,
  breakable,
  colback=palewarm,
  colframe=accenttwo,
  boxrule=0.6pt,
  arc=2pt,
  left=9pt,right=9pt,top=8pt,bottom=8pt,
  title={#1},
  fonttitle=\bfseries
}

\lstdefinestyle{cleanlisting}{
  basicstyle=\ttfamily\small,
  backgroundcolor=\color{codegray},
  frame=single,
  rulecolor=\color{rulegray},
  framerule=0.4pt,
  xleftmargin=0.5em,
  xrightmargin=0.5em,
  aboveskip=0.8em,
  belowskip=0.8em,
  columns=fullflexible,
  keepspaces=true,
  showstringspaces=false,
  breaklines=true,
  keywordstyle=\bfseries\color{accent},
  commentstyle=\itshape\color{accenttwo}
}
\lstset{style=cleanlisting}

\newcommand{\Lambert}{W}
\newcommand{\rLambert}{\mathcal W}
\newcommand{\dd}{\,\mathrm d}
\newcommand{\e}{\mathrm e}
\newcommand{\D}{\frac{\mathrm d}{\mathrm dz}}
\newcommand{\deltaop}{\delta}
\newcommand{\bigO}{\mathcal O}
\newcommand{\littleo}{o}
\newcommand{\fall}[2]{\left(#1\right)_{\!\underline{#2}}}
\newcommand{\bracks}[2]{\genfrac{[}{]}{0pt}{}{#1}{#2}}
\newcommand{\Trunc}{\operatorname{Trunc}}
\newcommand{\Res}{\operatorname{Res}}
\newcommand{\val}{\operatorname{val}}

\begin{document}
\hypersetup{pageanchor=false}

\begin{titlepage}
\thispagestyle{empty}
\vspace*{1.0in}

{\color{accent}\rule{\textwidth}{1.1pt}}\\[1.2em]
{\Huge\bfseries\color{ink}
Reversing \(x+W(x)\)}\\[0.65em]
{\fontsize{17}{21}\selectfont\color{accent}
Exact reduction, asymptotic sectors, and a calculus for logarithmic transseries}\\[1.2em]
{\color{accent}\rule{\textwidth}{1.1pt}}

\vspace{1.2in}

\begin{mainbox}[Central formula]
Let \(W\) be the principal real Lambert function and let
\[
  f(x)=x+W(x),\qquad x\ge 0.
\]
Writing \(g=f^{-1}\), \(w=W(z)\), and
\(
Lambert_r\) for the inverse of \(u\mapsto u\e^u+ru\), one has
\[
  \boxed{\;g(z)=z-\rLambert_1(z)\;}
\]
and, as \(z\to\infty\),
\[
  \boxed{\;
  g(z)=z-w+\sum_{k=1}^{N}\frac{A_k(w)}{z^k}
       +\bigO_N\!\left(\left(\frac{w}{z}\right)^{N+1}\right),
  \qquad
  A_k(w)=\frac{(-1)^{k+1}}{(k+1)!}
  \mathcal L_{k-1}\cdots\mathcal L_0\bigl(w^{k+1}\bigr),
  \;}
\]
where
\[
  \mathcal L_m=\frac{w}{1+w}\frac{\mathrm d}{\mathrm dw}-m.
\]
The small parameter is \(w/z=\e^{-w}\sim (\log z)/z\).  This sector is
beyond every inverse power of \(\log z\).
\end{mainbox}

\vfill
\begin{tabularx}{\textwidth}{@{}Xr@{}}
\textbf{Prepared as a self-contained research article} & 2 September 2026\\
\textit{Exact symbolic coefficients, proofs, algorithms, and numerical checks} &
\end{tabularx}
\end{titlepage}
\hypersetup{pageanchor=true}
\pagenumbering{roman}

\begin{abstract}
We determine the inverse of \(f(x)=x+W(x)\) on the nonnegative real axis and
construct its complete asymptotic description at infinity.  An elementary
change of variable reduces the inverse exactly to the \(r\)-Lambert function:
if \(\rLambert_r\) denotes the inverse of \(u\e^u+ru\), then
\(f^{-1}(z)=z-\rLambert_1(z)\).  Relative to the ordinary Lambert value
\(w=W(z)\), the difference is governed by the exponentially small parameter
\(\eta=w/z=\e^{-w}\).  We obtain a convergent local expansion in \(\eta\), a
uniform fixed-order remainder estimate, and an explicit differential-operator
formula for every coefficient.  The first coefficients are rational functions
of \(w\), beginning with
\[
 \frac{w^2}{1+w},\qquad
 \frac{w^3(w^2-2)}{2(1+w)^3},\qquad
 \frac{w^4(2w-1)(w^3-6w-6)}{6(1+w)^5}.
\]
Combining these sectors with the classical logarithmic expansion of \(W(z)\)
produces a nested transseries in \(z^{-1}\), \(\log z\), and
\(\log\log z\).  We then formulate a general reversal calculus for maps
\(F(z)=z+\Phi(z)\) whose perturbation is a logarithmic transseries.  In the
filtered differential algebra with \(q=z^{-1}\) and
\(\delta=z\,\mathrm d/\mathrm dz\), the identity
\(
\mathrm D(q^m a)=q^{m+1}(\delta-m)a
\)
turns Lagrange--B\"urmann inversion into a triangular, finitely computable
algorithm.  Fixed-point iteration, Newton iteration, truncation rules, and
residual certificates are developed in parallel.  The result is both an
explicit solution of the stated inverse problem and a reusable method for
series reversal across a broad class of polynomial-logarithmic transseries.
\end{abstract}

\tableofcontents
\clearpage
\pagenumbering{arabic}

\section{Problem, notation, and principal conclusions}

Throughout, \(W=W_0\) denotes the principal real Lambert function, characterized
on \([0,\infty)\) by \cite{Corless1996,DLMF}
\begin{equation}\label{eq:Wdef}
  W(x)\e^{W(x)}=x,\qquad W(x)\ge 0.
\end{equation}
We study
\begin{equation}\label{eq:fdef}
  f(x)=x+W(x),\qquad x\ge 0,
\end{equation}
and its inverse
\[
  g(z)=f^{-1}(z),\qquad z\ge 0.
\]
For the asymptotic analysis we use
\begin{equation}\label{eq:notation}
  w=W(z),\qquad L=\log z,
  \qquad \ell=\log L,
  \qquad q=z^{-1},
  \qquad \eta=\frac{w}{z}=\e^{-w}.
\end{equation}
The identities \(z=w\e^w\) and \(\eta=\e^{-w}\) will repeatedly convert
ordinary derivatives into operations on rational functions of \(w\).

There are two distinct asymptotic layers.

\begin{enumerate}[label=\textbf{\arabic*.},leftmargin=2.2em]
\item The familiar logarithmic layer is the Poincar\'e expansion
\[
 W(z)\sim L-\ell+\frac{\ell}{L}
 +\frac{\ell(\ell-2)}{2L^2}+\cdots.
\]
Consequently, the entire algebraic-logarithmic expansion of \(g\) begins as
\(z-W(z)\).

\item The difference between \(g(z)\) and \(z-W(z)\) is not another inverse
power of \(L\).  It is of order
\[
 \frac{w^2}{z(1+w)}\sim\frac{\log z}{z},
\]
which is smaller than \(L^{-M}\) for every fixed \(M\).  It therefore belongs
to a new, exponentially small transseries sector.
\end{enumerate}

The following statement summarizes the exact and asymptotic answers.  The organization by algebra, differentiation, composition, and reversion is also the organizing perspective of the linked ProveIt research-draft directory \cite{ProveItDrafts}.

\begin{theorem}[Main inversion theorem]\label{thm:main}
Let \(\rLambert_r\) denote the real inverse, near the nonnegative axis, of
\begin{equation}\label{eq:rLambertDef}
  u\longmapsto u\e^u+ru.
\end{equation}
Then
\begin{equation}\label{eq:exactinverse}
  g(z)=z-\rLambert_1(z),\qquad z\ge0.
\end{equation}
If \(w=W(z)\), then for every fixed integer \(N\ge0\),
\begin{equation}\label{eq:mainexpansion}
  g(z)=z-w+\sum_{k=1}^{N}\frac{A_k(w)}{z^k}
  +\bigO_N\!\left(\left(\frac{w}{z}\right)^{N+1}\right),
  \qquad z\to\infty,
\end{equation}
where
\begin{equation}\label{eq:Akoperator}
  A_k(w)=\frac{(-1)^{k+1}}{(k+1)!}
  \mathcal L_{k-1}\mathcal L_{k-2}\cdots\mathcal L_0
  \bigl(w^{k+1}\bigr),
  \qquad
  \mathcal L_m=\frac{w}{1+w}\frac{\mathrm d}{\mathrm dw}-m.
\end{equation}
Moreover,
\begin{equation}\label{eq:Akleading}
  A_k(w)=\frac{w^k}{k}+\bigO(w^{k-1})\qquad(w\to\infty).
\end{equation}
\end{theorem}

The exact reduction in \cref{eq:exactinverse} is simple, but it does not by
itself display the hierarchy of logarithmic and exponentially small terms.
The remainder of the article develops that hierarchy and then extracts a
general reversal mechanism that no longer depends specifically on \(W\).

\section{Exact inversion through the \texorpdfstring{\(r\)}{r}-Lambert function}

\subsection{Existence and monotonicity}

For \(x>0\), differentiation of \cref{eq:Wdef} gives
\begin{equation}\label{eq:Wprime}
  W'(x)=\frac{W(x)}{x(1+W(x))}
       =\frac{1}{\e^{W(x)}(1+W(x))}>0,
\end{equation}
with the continuous limiting value \(W'(0)=1\).  Hence
\[
  f'(x)=1+W'(x)>1
\]
for \(x\ge0\).  The map \(f\) is therefore strictly increasing from
\([0,\infty)\) onto \([0,\infty)\), so its inverse \(g\) exists globally on
that interval.

The same monotonicity gives an especially useful a posteriori estimate.  If
\(\widetilde g(z)\ge0\) is any approximation and
\begin{equation}\label{eq:residualdef}
  R(z)=f(\widetilde g(z))-z,
\end{equation}
then the mean-value theorem yields
\begin{equation}\label{eq:residualbound}
  \bigl|\widetilde g(z)-g(z)\bigr|\le |R(z)|.
\end{equation}
Thus a small residual is already a rigorous absolute-error certificate.

\subsection{The exact change of variable}

Set
\[
  u=W(g(z)).
\]
Then \(g(z)=u\e^u\), while \(z=g(z)+W(g(z))\).  Consequently
\begin{equation}\label{eq:uequation}
  z=u\e^u+u.
\end{equation}
On \(u\ge0\), the right-hand side has derivative
\(\e^u(1+u)+1>0\), so \cref{eq:uequation} has exactly one nonnegative
solution.  With the notation \cref{eq:rLambertDef}, that solution is
\(u=\rLambert_1(z)\).  Since \(g=u\e^u=z-u\), we obtain
\[
  g(z)=z-\rLambert_1(z).
\]
This proves the exact part of \cref{thm:main}.

\begin{remark}[Notation]\label{rem:notation}
The notation follows the generalized Lambert function introduced for the inverse of \(u\e^u+ru\) in \cite{MezoBaricz2017}.  The subscript in \(\rLambert_r\) is a deformation parameter, not a branch
index of the ordinary Lambert function.  We reserve \(W\) for the principal
ordinary branch and use calligraphic \(\rLambert\) precisely to avoid that
ambiguity.
\end{remark}

\subsection{A one-parameter family}

The same substitution solves a larger family without additional work.
For
\begin{equation}\label{eq:fr}
  f_r(x)=x+rW(x),\qquad r\ge0,
\end{equation}
putting \(u=W(f_r^{-1}(z))\) gives
\[
  z=u\e^u+ru,
  \qquad
  f_r^{-1}(z)=u\e^u=z-ru.
\]
Therefore
\begin{equation}\label{eq:frinverse}
  f_r^{-1}(z)=z-r\rLambert_r(z).
\end{equation}
The original problem is the case \(r=1\).  We will retain \(r\) in several
intermediate formulas because it reveals the organization of the transseries
coefficients.

\section{The hidden small parameter and an exact difference equation}

The ordinary value \(w=W(z)\) solves \(w\e^w=z\), whereas
\(u=\rLambert_r(z)\) solves \(u\e^u+ru=z\).  Define
\begin{equation}\label{eq:ddef}
  d=w-u,
  \qquad
  \eta=\frac{w}{z}=\e^{-w}.
\end{equation}
Substituting \(u=w-d\) into the equation for \(u\), and dividing by
\(z=w\e^w\), gives the exact identity
\begin{equation}\label{eq:differenceeq}
  \boxed{\quad
  \left(1-\frac{d}{w}\right)\bigl(\e^{-d}+r\eta\bigr)=1.
  \quad}
\end{equation}
For fixed large \(w\), this is an implicit analytic equation for \(d\) in the
small parameter \(r\eta\).

The scale \(\eta\) is exponentially small in \(w\), but only algebraically
small in \(z\):
\begin{equation}\label{eq:etasize}
  \eta=\e^{-w}=\frac{w}{z}\sim\frac{\log z}{z}.
\end{equation}
Since \(z^{-1}=\e^{-L}\), one has
\begin{equation}\label{eq:beyondorders}
  \frac{L}{z}=o(L^{-M})
  \qquad\text{for every fixed }M>0.
\end{equation}
This is the precise reason that the correction cannot be detected by carrying
the ordinary expansion of \(W(z)\) to any finite number of inverse logarithmic
orders.

\subsection{Analytic expansion in \texorpdfstring{\(\eta\)}{eta}}

Write \(\xi=r\eta\) and define
\[
  H(d,\xi;w)=\left(1-\frac{d}{w}\right)(\e^{-d}+\xi)-1.
\]
At \((d,\xi)=(0,0)\),
\[
  H(0,0;w)=0,
  \qquad
  \frac{\partial H}{\partial d}(0,0;w)=-\left(1+\frac1w\right)\ne0.
\]
The analytic implicit-function theorem therefore gives a unique germ
\begin{equation}\label{eq:dseries}
  d(w,\xi)=\sum_{k\ge1}b_k(w)\xi^k.
\end{equation}
For \(w\ge1\), the germ and a sufficiently small disk in \(\xi\) may be
chosen uniformly: introduce \(t=1/w\in[0,1]\), observe that
\(H=(1-td)(\e^{-d}+\xi)-1\), and apply compactness to the nonvanishing
of \(H_d\) near \((d,\xi)=(0,0)\).

The first derivative at the origin is
\begin{equation}\label{eq:b1}
  b_1(w)=-\frac{H_\xi}{H_d}(0,0;w)=\frac{w}{1+w}.
\end{equation}
Because
\[
  f_r^{-1}(z)=z-r(w-d)=z-rw+r d
\]
and \(\xi=r w/z\), \cref{eq:dseries} becomes
\begin{equation}\label{eq:grA}
  f_r^{-1}(z)
  =z-rw+\sum_{k\ge1}r^{k+1}\frac{A_k(w)}{z^k},
  \qquad
  A_k(w)=w^k b_k(w).
\end{equation}
For \(r=1\), this is the desired expansion of \(g\).

\begin{theorem}[Uniform fixed-order remainder]\label{thm:uniformremainder}
For each fixed \(N\ge0\) and each bounded interval \(0\le r\le R\),
\begin{equation}\label{eq:generalremainder}
  f_r^{-1}(z)=z-rw+
  \sum_{k=1}^{N}r^{k+1}\frac{A_k(w)}{z^k}
  +\bigO_{N,R}\!\left(\left(\frac{w}{z}\right)^{N+1}\right),
  \qquad z\to\infty.
\end{equation}
\end{theorem}

\begin{proof}
The analytic solution \(d(w,\xi)\) has a Taylor remainder bounded uniformly
for \(w\ge1\) and \(|\xi|\le\rho\), for some fixed \(\rho>0\).  Since
\(\xi=r w/z\to0\) uniformly for \(0\le r\le R\), truncating
\cref{eq:dseries} after \(N\) terms gives
\(d-\sum_{k=1}^N b_k\xi^k=\bigO_{N,R}(\xi^{N+1})\).  Multiplication by the
outer factor \(r\) in \(f_r^{-1}=z-rw+rd\) only changes the constant when
\(r\) is bounded.  Substitution of \(A_k=w^k b_k\) proves the claim.
\end{proof}

\subsection{Leading coefficient at every sector}

The difference equation also identifies the dominant term in every
\(z^{-k}\) sector.  Let \(w\to\infty\) in \cref{eq:differenceeq} while keeping
\(\xi\) fixed near zero.  The limiting equation is
\[
  \e^{-d}+\xi=1,
\]
so
\begin{equation}\label{eq:resummedleading}
  d\longrightarrow-\log(1-\xi)
  =\sum_{k\ge1}\frac{\xi^k}{k}.
\end{equation}
Consequently
\[
  b_k(w)\longrightarrow\frac1k,
  \qquad
  A_k(w)=\frac{w^k}{k}+\bigO(w^{k-1}).
\]
The leading pieces of all exponentially small sectors are therefore the
Taylor coefficients of a single elementary function, \(-\log(1-\eta)\).
This does not replace the full coefficient functions, but it is a useful
large-\(w\) resummation principle.

\section{Lagrange--B\"urmann reversal for near-identity maps}

The coefficient functions in \cref{eq:grA} can be obtained directly from a
general inversion identity.  This section develops the identity first in
ordinary differential form and then in the filtered differential algebra
adapted to logarithmic transseries.

\subsection{The universal derivative formula}

Let
\begin{equation}\label{eq:nearidentity}
  F(x)=x+\Phi(x)
\end{equation}
with \(\Phi(x)=o(x)\) in the asymptotic regime of interest.  Introduce a
bookkeeping parameter \(\varepsilon\) and write
\[
  G_\varepsilon(z)=z+U_\varepsilon(z),
  \qquad
  U_\varepsilon=-\varepsilon\Phi(z+U_\varepsilon).
\]
Lagrange--B\"urmann inversion \cite{FlajoletSedgewick,Henrici} gives the formal identity
\begin{equation}\label{eq:LB}
  G_\varepsilon(z)=z+
  \sum_{n\ge1}\frac{(-\varepsilon)^n}{n!}
  \frac{\mathrm d^{n-1}}{\mathrm dz^{n-1}}
  \bigl(\Phi(z)^n\bigr).
\end{equation}
Whenever the right-hand side is summable in the relevant asymptotic filtration,
we set \(\varepsilon=1\) and obtain
\begin{equation}\label{eq:LB1}
  \boxed{\quad
  F^{-1}(z)=z+
  \sum_{n\ge1}\frac{(-1)^n}{n!}
  \frac{\mathrm d^{n-1}}{\mathrm dz^{n-1}}
  \bigl(\Phi(z)^n\bigr).
  \quad}
\end{equation}
The first four nontrivial blocks are
\begin{align}
 F^{-1}(z)={}&z-\Phi
 +\Phi\Phi'
 -\Phi(\Phi')^2-\frac12\Phi^2\Phi'' \notag\\
 &+\Phi(\Phi')^3
 +\frac32\Phi^2\Phi'\Phi''
 +\frac16\Phi^3\Phi'''
 +\cdots.\label{eq:universalfirstterms}
\end{align}
For logarithmic \(\Phi\), each ordinary derivative introduces one factor of
\(z^{-1}\), so the displayed lines are successive transseries sectors.

\subsection{The coefficient differential algebra}

Set
\begin{equation}\label{eq:qdelta}
  q=z^{-1},
  \qquad
  \delta=z\frac{\mathrm d}{\mathrm dz}.
\end{equation}
Let \(\mathcal B\) be a differential algebra of slowly varying coefficient
functions, stable under \(\delta\).  Typical elements are finite or formal
combinations of
\[
  L=\log z,\qquad \ell=\log L,
  \qquad \log\ell,\qquad L^{-1},\qquad \ell^{-1},\ldots,
\]
with a chosen depth and a chosen asymptotic ordering; this is a deliberately elementary fragment of the differential-algebra viewpoint used for transseries \cite{vanDerHoeven}.  Work in the filtered
ring \(\mathcal B[[q]]\).  For \(a\in\mathcal B\), direct differentiation gives
\begin{equation}\label{eq:Dblock}
  \frac{\mathrm d}{\mathrm dz}\bigl(q^m a\bigr)
  =q^{m+1}(\delta-m)a.
\end{equation}
Thus differentiation raises the \(q\)-valuation by at least one.  In
particular, for \(a\in\mathcal B\),
\begin{equation}\label{eq:Dkblock}
  \frac{\mathrm d^k a}{\mathrm dz^k}
  =q^k\fall{\delta}{k}a,
  \qquad
  \fall{\delta}{k}
  =\delta(\delta-1)\cdots(\delta-k+1).
\end{equation}
The operator on the right is a falling factorial in \(\delta\).

\begin{theorem}[Pure logarithmic perturbation]\label{thm:purelogreverse}
Suppose \(\Phi\in\mathcal B\) and \(q\Phi\prec1\).  Then the formal inverse of
\(F(z)=z+\Phi(z)\) is
\begin{equation}\label{eq:purelogformula}
  F^{-1}(z)=z-\Phi(z)
  +\sum_{k\ge1}q^k\frac{(-1)^{k+1}}{(k+1)!}
  \fall{\delta}{k}\bigl(\Phi^{k+1}\bigr).
\end{equation}
For any prescribed \(q\)-order \(N\), only the terms \(k\le N\) can
contribute.
\end{theorem}

\begin{proof}
Put \(n=k+1\) in \cref{eq:LB1} and use \cref{eq:Dkblock}.  The \(k\)-th
summand has \(q\)-valuation at least \(k\), so coefficient extraction through
\(q^N\) is finite and triangular.
\end{proof}

\begin{remark}[General \(q\)-dependent perturbations]\label{rem:qdependent}
If
\(
\Phi=\sum_{j\ge0}q^j\phi_j
\)
with \(\phi_j\in\mathcal B\), the original formula \cref{eq:LB1} still
applies.  One simply evaluates every derivative by the block rule
\cref{eq:Dblock} and truncates after each operation.  Because the \(n\)-th
Lagrange term begins at \(q^{n-1}\), terms with \(n>N+1\) cannot affect the
inverse through \(q^N\).
\end{remark}

\section{Specialization to \texorpdfstring{\(\Phi=W\)}{Phi = W}}

For \(w=W(z)\), \cref{eq:Wprime} becomes
\begin{equation}\label{eq:deltaw}
  \delta w=zW'(z)=\frac{w}{1+w}.
\end{equation}
If \(R=R(w)\), then
\begin{equation}\label{eq:DzR}
  \frac{\mathrm d}{\mathrm dz}\bigl(z^{-m}R(w)\bigr)
  =z^{-m-1}\left(
    \frac{w}{1+w}\frac{\mathrm dR}{\mathrm dw}-mR
  \right).
\end{equation}
This is precisely the operator \(\mathcal L_m\) in \cref{eq:Akoperator}.
Applying \cref{thm:purelogreverse} to \(F(x)=x+W(x)\) proves the coefficient
formula in \cref{thm:main}.

\subsection{First coefficient functions}

The first four rational functions are
\begin{align}
 A_1(w)&=\frac{w^2}{1+w},\label{eq:A1}\\[0.25em]
 A_2(w)&=\frac{w^3(w^2-2)}{2(1+w)^3},\label{eq:A2}\\[0.25em]
 A_3(w)&=\frac{w^4(2w-1)(w^3-6w-6)}{6(1+w)^5},\label{eq:A3}\\[0.25em]
 A_4(w)&=\frac{w^5}{24(1+w)^7}
 \bigl(6w^6-8w^5-69w^4-40w^3
       +96w^2+72w-24\bigr).\label{eq:A4}
\end{align}
The fifth and sixth coefficients, useful for symbolic and numerical checking,
are recorded in \cref{app:higherA}.  In general, induction in
\cref{eq:Akoperator} shows that
\begin{equation}\label{eq:Akstructure}
  A_k(w)\in
  \frac{w^{k+1}}{(1+w)^{2k-1}}\,\mathbb Q[w],
\end{equation}
where the remaining numerator polynomial has degree at most \(2k-2\).

Thus the inverse begins
\begin{align}
 g(z)={}&z-w
 +\frac{1}{z}\frac{w^2}{1+w}
 +\frac{1}{z^2}\frac{w^3(w^2-2)}{2(1+w)^3}\notag\\
 &+\frac{1}{z^3}
   \frac{w^4(2w-1)(w^3-6w-6)}{6(1+w)^5}
 +\bigO\!\left(\left(\frac{w}{z}\right)^4\right).
 \label{eq:firstfour}
\end{align}
Because the coefficient functions are evaluated at the exact \(w=W(z)\),
\cref{eq:firstfour} is already a highly effective asymptotic formula without
expanding any logarithms.

\subsection{Large-\texorpdfstring{\(w\)}{w} form of the coefficients}

Expanding the rational functions at \(w=\infty\) gives
\begin{align}
 A_1(w)={}&w-1+\frac1w-\frac1{w^2}+\frac1{w^3}
 +\bigO(w^{-4}),\label{eq:A1large}\\
 A_2(w)={}&\frac{w^2}{2}-\frac{3w}{2}+2-\frac2w
 +\frac{3}{2w^2}-\frac{1}{2w^3}+\bigO(w^{-4}),\label{eq:A2large}\\
 A_3(w)={}&\frac{w^3}{3}-\frac{11w^2}{6}
 +\frac{23w}{6}-\frac{31}{6}+\frac{31}{6w}
 -\frac{11}{3w^2}+\bigO(w^{-3}),\label{eq:A3large}\\
 A_4(w)={}&\frac{w^4}{4}-\frac{25w^3}{12}
 +\frac{155w^2}{24}-\frac{95w}{8}+\frac{47}{3}
 +\bigO(w^{-1}).\label{eq:A4large}
\end{align}
The leading terms agree with \cref{eq:Akleading}.  Retaining only those
leading terms in every sector produces
\[
  g(z)-(z-w)\sim
  \sum_{k\ge1}\frac1k\left(\frac{w}{z}\right)^k
  =-\log\left(1-\frac{w}{z}\right),
\]
which is the sector-leading resummation already found from the implicit
difference equation.

\section{The classical logarithmic expansion of \texorpdfstring{\(W\)}{W}}

To express the answer solely through elementary iterated logarithms, write
\[
  L=\log z,
  \qquad
  \ell=\log L.
\]
Let \(\bracks{n}{j}\) denote the unsigned Stirling number of the first kind.
Define
\begin{equation}\label{eq:Pndef}
  P_n(\ell)=\sum_{m=1}^{n}
  \frac{(-1)^{n-m}}{m!}
  \bracks{n}{n-m+1}\ell^m.
\end{equation}
Then the standard asymptotic expansion \cite{Corless1996,DLMF} is
\begin{equation}\label{eq:Wlogseries}
  W(z)\sim L-\ell+\sum_{n\ge1}\frac{P_n(\ell)}{L^n}.
\end{equation}
The first polynomials are
\begin{align}
 P_1(\ell)&=\ell,\notag\\
 P_2(\ell)&=\frac{\ell(\ell-2)}2,\notag\\
 P_3(\ell)&=\frac{\ell(2\ell^2-9\ell+6)}6,\notag\\
 P_4(\ell)&=\frac{\ell(3\ell^3-22\ell^2+36\ell-12)}{12},\notag\\
 P_5(\ell)&=\frac{\ell(12\ell^4-125\ell^3+350\ell^2-300\ell+60)}{60}.
 \label{eq:Pfirst}
\end{align}
Further polynomials through \(P_8\) are listed in \cref{app:Ppolys}.

For each fixed \(M\), the expansion has the usual Poincar\'e interpretation \cite{deBruijn}:
\begin{equation}\label{eq:Wlogremainder}
 W(z)=L-\ell+\sum_{n=1}^{M}\frac{P_n(\ell)}{L^n}
 +\bigO\!\left(\frac{\ell^{M+1}}{L^{M+1}}\right).
\end{equation}
Since the first beyond-all-orders correction is \(\bigO(L/z)\), the purely
logarithmic Poincar\'e expansion of the inverse is
\begin{equation}\label{eq:goutergeneral}
  g(z)\sim z-L+\ell-\sum_{n\ge1}\frac{P_n(\ell)}{L^n}.
\end{equation}
Explicitly,
\begin{align}
 g(z)\sim{}&z-L+\ell-\frac{\ell}{L}
 +\frac{\ell(2-\ell)}{2L^2}
 -\frac{\ell(2\ell^2-9\ell+6)}{6L^3}\notag\\
 &+\frac{\ell(12-36\ell+22\ell^2-3\ell^3)}{12L^4}
 -\frac{P_5(\ell)}{L^5}-\cdots.
 \label{eq:gouterexplicit}
\end{align}
At every fixed inverse-logarithmic order this is correct, but it cannot reveal
the term \(A_1(w)/z\sim L/z\).  That term is smaller than the remainder of
\emph{every} finite truncation in \cref{eq:gouterexplicit}.  This is a standard
instance of information that is absent from a single Poincar\'e scale but present
in a transseries.

\section{The full nested logarithmic transseries}

Combining \cref{eq:mainexpansion} and \cref{eq:Wlogseries} produces the
sector decomposition
\begin{equation}\label{eq:fulltransseriesform}
  g(z)\sim z+\mathcal G_0(L,\ell)
  +\sum_{k\ge1}z^{-k}\mathcal G_k(L,\ell),
\end{equation}
where
\begin{equation}\label{eq:Gsectors}
  \mathcal G_0=-W_{\mathrm{log}},
  \qquad
  \mathcal G_k=A_k(W_{\mathrm{log}}),\quad k\ge1,
\end{equation}
with \(W_{\mathrm{log}}\) denoting the formal series on the right of
\cref{eq:Wlogseries}.

The ordering must be read sector by sector.  All inverse powers of \(L\) in
\(\mathcal G_0\) precede the first \(z^{-1}\) term; all inverse powers of
\(L\) in the \(z^{-1}\) sector precede the first \(z^{-2}\) term; and so on.
Symbolically,
\[
  L^{-M}\gg z^{-1}L^J
  \qquad\text{for every fixed }M,J.
\]
This lexicographic ordering is essential.  Appending a \(z^{-1}\) correction
after only a few terms of the outer expansion does not improve a numerical
approximation unless the neglected outer remainder has first been made smaller
than the new sector.

\subsection{First three exponentially small sectors}

Straight substitution and re-expansion give
\begin{align}
\mathcal G_1(L,\ell)={}&
 L-\ell-1+\frac{\ell+1}{L}
 +\frac{\ell^2-2}{2L^2}
 +\frac{2\ell^3-3\ell^2-12\ell+6}{6L^3}
 +\bigO\!\left(\frac{\ell^4}{L^4}\right),
 \label{eq:G1exp}\\[0.4em]
\mathcal G_2(L,\ell)={}&
 \frac{L^2}{2}-\frac{2\ell+3}{2}L
 +\frac{(\ell+1)(\ell+4)}2
 -\frac{(\ell+1)(\ell+4)}{2L}
 +\bigO\!\left(\frac{\ell^3}{L^2}\right),
 \label{eq:G2exp}\\[0.4em]
\mathcal G_3(L,\ell)={}&
 \frac{L^3}{3}-\frac{6\ell+11}{6}L^2
 +\frac{6\ell^2+28\ell+23}{6}L\notag\\
 &-\frac{2\ell^3+20\ell^2+51\ell+31}{6}
 +\bigO\!\left(\frac{\ell^3}{L}\right).
 \label{eq:G3exp}
\end{align}
Thus, in fully elementary form,
\begin{align}
 g(z)\sim{}&z-L+\ell-\frac{\ell}{L}
 +\frac{\ell(2-\ell)}{2L^2}+\cdots\notag\\
 &+\frac1z\left[
 L-\ell-1+\frac{\ell+1}{L}
 +\frac{\ell^2-2}{2L^2}+\cdots\right]\notag\\
 &+\frac1{z^2}\left[
 \frac{L^2}{2}-\frac{2\ell+3}{2}L
 +\frac{(\ell+1)(\ell+4)}2+\cdots\right]\notag\\
 &+\frac1{z^3}\left[
 \frac{L^3}{3}-\frac{6\ell+11}{6}L^2
 +\frac{6\ell^2+28\ell+23}{6}L+\cdots\right]+\cdots.
 \label{eq:fullfirstsectors}
\end{align}
The leading term in the \(k\)-th small sector is
\begin{equation}\label{eq:sectorleading}
  \frac{L^k}{kz^k}.
\end{equation}
Subleading powers of \(L\) and polynomials in \(\ell\) are produced
mechanically by \cref{eq:Akoperator} and \cref{eq:Pndef}.

\subsection{A practical two-level error budget}

There are two independent truncation decisions.

\begin{enumerate}[label=(\alph*),leftmargin=2em]
\item Keeping exact \(w=W(z)\) and truncating after the \(z^{-N}\) sector has
remainder \(\bigO((w/z)^{N+1})\) by \cref{thm:uniformremainder}.

\item Replacing a coefficient \(A_k(w)\) by its logarithmic expansion through
\(J\) subleading inverse powers of \(L\) introduces, at that sector, a typical
error
\[
  z^{-k}\bigO\bigl(L^{k-J-1}\ell^{J+1}\bigr),
\]
with the evident modification for the outer sector \(k=0\).
\end{enumerate}

For numerical work, the most efficient representation is usually the exact-
\(W\) sector expansion \cref{eq:mainexpansion}.  The fully elementary nested
form \cref{eq:fullfirstsectors} is most useful for structural analysis,
comparison of monomials, and integration into a formal transseries system.

\section{A general reversal method for logarithmic transseries}

We now detach the method from the special function \(W\).  Consider
\begin{equation}\label{eq:generalF}
  F(z)=z+\Phi(z),
  \qquad
  \Phi\in\mathcal B[[q]],
  \qquad q=z^{-1},
\end{equation}
where \(\mathcal B\) is a chosen logarithmic coefficient algebra stable under
\(\delta=z\,\mathrm d/\mathrm dz\).  Assume that \(q\Phi\prec1\), so the map
is asymptotically near the identity.

\subsection{Direct Lagrange algorithm}

The most direct algorithm is a literal implementation of \cref{eq:LB1}, with
all products and derivatives truncated at the requested \(q\)-order.

\begin{methodbox}[Algorithm 1: Lagrange--B\"urmann reversal through \(q^N\)]
\begin{enumerate}[leftmargin=2em]
\item Represent \(\Phi\) as a truncated block series
\(\sum_{m=0}^{N}q^m\phi_m\), with \(\phi_m\in\mathcal B\).
\item Initialize \(G\leftarrow z\).
\item For \(n=1,\ldots,N+1\):
  form \(H\leftarrow\Phi^n\), truncate at \(q^N\), apply \(n-1\)
  derivatives using
  \[
    D(q^m a)=q^{m+1}(\delta-m)a,
  \]
  truncating after each derivative, and add
  \((-1)^nH/n!\) to \(G\).
\item Return \(G\) truncated through \(q^N\).
\end{enumerate}
Only finitely many coefficient operations are performed, because the
\(n\)-th summand begins at \(q^{n-1}\).
\end{methodbox}

A compact pseudocode version is as follows.

\begin{lstlisting}[language=Python,caption={Filtered Lagrange reversal.}]
def reverse_near_identity(phi, N):
    # phi is a block series sum(q**m * a[m]), truncated at q**N.
    G = z
    for n in range(1, N + 2):
        H = truncate(phi**n, q_order=N)
        for _ in range(n - 1):
            H = truncate(D_block(H), q_order=N)
        G += (-1)**n * H / factorial(n)
    return truncate(G, q_order=N)

# On one block q**m * a:
# D_block(q**m * a) = q**(m+1) * (delta(a) - m*a)
\end{lstlisting}

\subsection{Triangular fixed-point recursion}

Writing the inverse as \(G=z+U\), one may instead solve
\begin{equation}\label{eq:fixedpointU}
  U=-\Phi(z+U).
\end{equation}
Taylor composition gives
\begin{equation}\label{eq:TaylorPhi}
  \Phi(z+U)=\sum_{j\ge0}\frac{U^j}{j!}\Phi^{(j)}(z).
\end{equation}
Each derivative raises \(q\)-valuation, so the coefficient of \(q^m\) on the
right depends only on already known lower blocks and the current unknown block
linearly.  This produces a triangular coefficient recursion.  It is often the
best method when the coefficient algebra already has a robust composition
engine.

\subsection{Newton iteration in the filtered algebra}

A third method starts from an approximation \(G_0\), for example
\(G_0=z-\Phi(z)\), and iterates
\begin{equation}\label{eq:Newtoninverse}
  G_{j+1}=G_j-
  \frac{F(G_j)-z}{F'(G_j)}.
\end{equation}
Since \(F'=1+o(1)\) is a unit, Newton correction is well-defined in the
filtered ring.  Under the usual completeness assumptions it approximately
doubles the number of correct \(q\)-blocks at each step.  Newton iteration is
therefore preferable at high truncation order, while the Lagrange formula is
preferable for closed coefficient formulas and proofs.

\subsection{Certification and two-sided checks}

For a symbolic candidate \(G\), compute the residual
\begin{equation}\label{eq:formalresidual}
  \mathcal R=F\circ G-z.
\end{equation}
If \(\mathcal R=\bigO(q^{N+1})\), then \(G\) is a right inverse through the
requested order.  A two-sided check also computes \(G\circ F-z\).  In a
well-ordered near-identity class, uniqueness of the inverse makes one side
sufficient in principle, but checking both sides is a valuable implementation
test for composition, differentiation, and truncation conventions.

For real numerical evaluation, a derivative lower bound converts the residual
into an error bound.  If \(F'(x)\ge m>0\) on the interval joining
\(G(z)\) to the exact inverse, then
\begin{equation}\label{eq:generalresidualbound}
  \bigl|G(z)-F^{-1}(z)\bigr|
  \le \frac{|F(G(z))-z|}{m}.
\end{equation}
For \(F(x)=x+W(x)\), one may take \(m=1\), recovering
\cref{eq:residualbound}.

\subsection{Choosing a coefficient algebra and monomial order}

A practical transseries implementation must specify more than formulas.  It
must decide what constitutes a coefficient and how monomials are compared.
For the present problem a natural hierarchy is
\begin{equation}\label{eq:monomialhierarchy}
  z\;\gg\;L\;\gg\;\ell\;\gg\;1\;\gg\;L^{-1}\;\gg\cdots
  \gg\;z^{-1}L^J\;\gg\;z^{-2}L^K\;\gg\cdots
\end{equation}
for every fixed \(J,K\).  More generally, one can use a lexicographic pair
\((m,\nu)\), where \(m\) is the outer \(q\)-sector and \(\nu\) is the
internal logarithmic order.  The following rules prevent most implementation
errors.

\begin{enumerate}[label=\textbf{R\arabic*.},leftmargin=2.5em]
\item Truncate first by outer \(q\)-valuation, then by the internal order of
coefficient blocks.
\item Keep \(\delta\) closed on the coefficient algebra; introduce additional
iterated logarithms only when differentiation requires them.
\item Apply truncation after every multiplication, derivative, and composition.
\item Record a target remainder ideal, not merely a number of printed terms.
\item Verify the final answer by composition in the same ordering convention.
\end{enumerate}

\subsection{A generic logarithmic perturbation}

Suppose, for example,
\begin{equation}\label{eq:genericPhi}
  \Phi(z)=aL+b\ell+c,
  \qquad L=\log z,
  \qquad \ell=\log L.
\end{equation}
Then
\[
  \delta\Phi=a+\frac{b}{L},
  \qquad
  \delta^2\Phi=-\frac{b}{L^2},
\]
and \cref{eq:purelogformula} immediately gives
\begin{align}
 F^{-1}(z)={}&z-\Phi
 +q\,\Phi\,\delta\Phi\notag\\
 &-\frac{q^2}{6}(\delta-1)\delta(\Phi^3)
 +\frac{q^3}{24}(\delta-2)(\delta-1)\delta(\Phi^4)
 +\bigO(q^4\mathcal B).
 \label{eq:genericreverse}
\end{align}
No ad hoc re-expansion of \(\log(z+U)\) or \(\log\log(z+U)\) is required.
The falling-factorial operators encode all chain-rule and power-of-\(z\)
bookkeeping automatically.  This is the principal advantage of the
\((q,\delta)\) formulation.

\section{Proof of the operator coefficient formula}

Although \cref{eq:Akoperator} follows quickly from Lagrange inversion, it is
useful to display the mechanism in full.

For \(F(x)=x+W(x)\), \cref{eq:LB1} reads
\begin{equation}\label{eq:WLB}
  g(z)=z+\sum_{n\ge1}\frac{(-1)^n}{n!}
  \frac{\mathrm d^{n-1}}{\mathrm dz^{n-1}}w^n.
\end{equation}
The \(n=1\) term is \(-w\).  Put \(n=k+1\) in every remaining term.  Starting
from a function \(R(w)\) with no explicit power of \(z^{-1}\), one derivative
uses \(\mathcal L_0\), the next uses \(\mathcal L_1\), and after \(k\)
derivatives one has
\begin{equation}\label{eq:operatorinduction}
  \frac{\mathrm d^k}{\mathrm dz^k}R(w)
  =z^{-k}\mathcal L_{k-1}\cdots\mathcal L_0R(w).
\end{equation}
This follows inductively from \cref{eq:DzR}.  Taking \(R(w)=w^{k+1}\) in
\cref{eq:operatorinduction} and substituting into \cref{eq:WLB} gives
\[
 \frac{(-1)^{k+1}}{(k+1)!}z^{-k}
 \mathcal L_{k-1}\cdots\mathcal L_0(w^{k+1}),
\]
which is exactly \(z^{-k}A_k(w)\).

The Lagrange series and the analytic \(\eta\)-series are the same expansion in
different coordinates.  Indeed, \(d=w-\rLambert_1(z)=g(z)-(z-w)\), and
\cref{eq:mainexpansion} gives
\[
 d=\sum_{k\ge1}\frac{A_k(w)}{z^k}
  =\sum_{k\ge1}\frac{A_k(w)}{w^k}\eta^k.
\]
Comparison with \cref{eq:dseries} shows
\begin{equation}\label{eq:bkAk}
  b_k(w)=\frac{A_k(w)}{w^k}.
\end{equation}
Thus the implicit-function proof supplies convergence and a remainder bound,
while Lagrange inversion supplies the closed differential-operator formula.
Together they give more than either method alone.

\section{Numerical verification}

The exact value can be computed without directly inverting \(f\): solve
\begin{equation}\label{eq:numericalu}
  u\e^u+u=z
\end{equation}
for its nonnegative root and set \(g=z-u\).  Let
\begin{equation}\label{eq:GN}
  G_N(z)=z-w+\sum_{k=1}^{N}\frac{A_k(w)}{z^k},
  \qquad w=W(z),
\end{equation}
with \(G_0=z-w\).  High-precision evaluation gives the absolute errors in
\cref{tab:errors}.

\begin{table}[htbp]
\centering
\caption{Absolute errors \(\lvert G_N(z)-g(z)\rvert\) using exact \(w=W(z)\).}
\label{tab:errors}
\small
\begin{tabular}{@{}rrrrrr@{}}
\toprule
\(z\) & \(N=0\) & \(N=1\) & \(N=2\) & \(N=3\) & \(N=4\)\\
\midrule
\(10\) & \(1.1202\times10^{-1}\) & \(1.0462\times10^{-3}\) & \(2.9897\times10^{-4}\) & \(2.3371\times10^{-5}\) & \(6.8889\times10^{-7}\)\\
\(100\) & \(2.6355\times10^{-2}\) & \(2.1862\times10^{-4}\) & \(9.4975\times10^{-7}\) & \(2.3486\times10^{-8}\) & \(7.9461\times10^{-10}\)\\
\(10^3\) & \(4.4172\times10^{-3}\) & \(7.5875\times10^{-6}\) & \(1.3534\times10^{-8}\) & \(1.7161\times10^{-11}\) & \(1.1027\times10^{-14}\)\\
\(10^6\) & \(1.0464\times10^{-5}\) & \(4.9551\times10^{-11}\) & \(2.9302\times10^{-16}\) & \(1.8408\times10^{-21}\) & \(1.1635\times10^{-26}\)\\
\bottomrule
\end{tabular}
\end{table}

Representative exact values are
\begin{align*}
 g(10)&=8.366493829844153615806835\ldots,\\
 g(100)&=96.64072495463040645894733\ldots,\\
 g(10^3)&=994.7548143481392806156898\ldots,\\
 g(10^6)&=999988.6166523780211244849\ldots.
\end{align*}
The error pattern is consistent with the scale
\((w/z)^{N+1}\).  At \(z=10^6\), for example, the zeroth approximation
\(z-W(z)\) misses a term of order \(10^{-5}\), while four correction sectors
reduce the error below \(10^{-26}\).

The numerical experiment also illustrates the distinction between the outer
and inner expansions.  Increasing the number of inverse-logarithmic terms in
\(W(z)\) drives a log-only approximation toward \(z-W(z)\), not toward the
exact inverse beyond the floor
\[
  g(z)-(z-W(z))=W(z)-\rLambert_1(z)
  \sim\frac{W(z)^2}{z(1+W(z))}.
\]
Crossing that floor requires the \(z^{-1}\) sector.

\section{Extensions and variants}

\subsection{The full \texorpdfstring{\(r\)}{r}-family}

Combining \cref{eq:frinverse} with the coefficient construction gives
\begin{equation}\label{eq:rfull}
  f_r^{-1}(z)=z-rw+
  \sum_{k\ge1}r^{k+1}\frac{A_k(w)}{z^k},
  \qquad w=W(z),
\end{equation}
with the fixed-order remainder stated in \cref{thm:uniformremainder}.  Equivalently,
\begin{equation}\label{eq:rLambertExpansion}
  \rLambert_r(z)=w-
  \sum_{k\ge1}r^k\frac{A_k(w)}{z^k}.
\end{equation}
Thus one sequence \(A_k\) simultaneously describes every bounded value of the
parameter \(r\).

\subsection{Affine-scaled Lambert perturbations}

Consider the three-parameter map
\begin{equation}\label{eq:scaledF}
  F(x)=\alpha x+\beta W(\gamma x),
  \qquad \alpha,\beta,\gamma>0.
\end{equation}
Putting \(u=W(\gamma x)\) gives
\[
  \frac{\gamma z}{\alpha}
  =u\e^u+\frac{\beta\gamma}{\alpha}u.
\]
Therefore
\begin{equation}\label{eq:scaledinverse}
  \boxed{\quad
  F^{-1}(z)=\frac{z-\beta\,
  \rLambert_{\beta\gamma/\alpha}(\gamma z/\alpha)}{\alpha}.
  \quad}
\end{equation}
The asymptotic expansion follows from \cref{eq:rLambertExpansion} after
replacing \(z\) by \(\gamma z/\alpha\).  This exact reduction is useful when
normalizing more complicated inverse problems before invoking the general
transseries machinery.

\subsection{Perturbations built from iterated logarithms}

The differential-algebra method does not require a special-function reduction.
For maps such as
\[
  F(x)=x+a\log x+b\log\log x
  +\frac{c_0+c_1\log\log x}{\log x}
  +\frac{d_0}{x}+\cdots,
\]
one places the slowly varying expressions in \(\mathcal B\), the explicit
powers of \(x^{-1}\) in the outer \(q\)-filtration, and applies either
Algorithm~1, the fixed-point recursion, or Newton iteration.  The same block
derivative \cref{eq:Dblock} guarantees triangularity.  The only additional
work is to define a consistent internal ordering for the chosen depth of
iterated logarithms.

\section{Implementation checklist and common pitfalls}

The following checklist summarizes the method in a form suitable for a
computer-algebra implementation or a semi-formal proof development.

\begin{methodbox}[A robust reversal workflow]
\begin{enumerate}[leftmargin=2em]
\item \textbf{Normalize the map.}  Rewrite it as \(F=z+\Phi\), possibly after
an affine change of variables.
\item \textbf{Choose the outer parameter.}  Usually \(q=z^{-1}\); verify that
\(q\Phi\prec1\).
\item \textbf{Choose the coefficient algebra.}  Fix the allowed logarithmic
depth and ensure closure under \(\delta=zD_z\).
\item \textbf{State the ordering.}  Use outer \(q\)-valuation first and
logarithmic order second.
\item \textbf{Generate coefficients.}  Use Lagrange--B\"urmann for closed
formulas, a triangular fixed point for transparent recursion, or Newton for
high order.
\item \textbf{Truncate after every operation.}  This keeps all computations
finite and prevents irrelevant high sectors from expanding combinatorially.
\item \textbf{Compose back.}  Check \(F\circ G-z\) to the target remainder
ideal; preferably also check \(G\circ F-z\).
\item \textbf{Certify numerically.}  Where a derivative lower bound is
available, convert the residual directly into an error estimate.
\end{enumerate}
\end{methodbox}

Several mistakes recur in logarithmic reversal problems.

\begin{enumerate}[label=\textbf{P\arabic*.},leftmargin=2.5em]
\item \textbf{Mixing scales.}  A finite \(1/L\) expansion cannot resolve a
term of size \(L/z\).  The latter requires a new sector.
\item \textbf{Differentiating coefficients as constants.}  The operator
\(D(q^m a)=q^{m+1}(\delta-m)a\) contains both \(\delta a\) and the derivative
of \(q^m\).
\item \textbf{Expanding too early.}  Keeping \(w=W(z)\) exact makes the sector
coefficients compact and numerically stable.  Expand in \(L,\ell\) only when
an elementary nested form is required.
\item \textbf{Using a one-dimensional truncation counter.}  A transseries
remainder is an ideal in a multi-scale ordering, not merely ``the next printed
term.''
\item \textbf{Ignoring branch data.}  The present real nonnegative problem has
unique branches, but complex or negative-domain variants require explicit
branch selection before composition and inversion.
\end{enumerate}

\section{Conclusions}

The inverse of \(x+W(x)\) admits both an exact special-function representation
and a transparent complete asymptotic structure.  The exact representation is
\[
  f^{-1}(z)=z-\rLambert_1(z).
\]
With \(w=W(z)\), the practically useful expansion is
\[
  f^{-1}(z)=z-w+\frac{w^2}{z(1+w)}
  +\frac{w^3(w^2-2)}{2z^2(1+w)^3}+\cdots,
\]
and truncation after \(z^{-N}\) has error
\(\bigO((w/z)^{N+1})\).  The parameter \(w/z=\e^{-w}\) identifies the
correction as beyond all orders of the conventional inverse-logarithmic
expansion.  Re-expansion of each rational coefficient in
\(L=\log z\) and \(\ell=\log L\) yields the full nested logarithmic
transseries.

The broader reversal method is encoded by two formulas:
\[
  F^{-1}(z)=z+
  \sum_{n\ge1}\frac{(-1)^n}{n!}D^{n-1}(\Phi^n),
  \qquad F=z+\Phi,
\]
and
\[
  D(q^m a)=q^{m+1}(\delta-m)a,
  \qquad q=z^{-1},\quad \delta=zD.
\]
Together they make coefficient extraction finite, triangular, and compatible
with arbitrary logarithmic depth.  The same structure supports direct
Lagrange inversion, fixed-point recursion, Newton acceleration, and residual
certification.  The Lambert example is therefore not only solvable in closed
form; it is a model case for systematic reversal of polynomial-logarithmic
transseries.

\appendix

\section{Higher rational coefficient functions}\label{app:higherA}

For compactness, write
\begin{equation}\label{eq:A5}
 A_5(w)=\frac{w^6}{120(1+w)^9}\,Q_5(w),
\end{equation}
where
\begin{align*}
 Q_5(w)={}&24w^8-58w^7-428w^6-151w^5
 +1320w^4+1380w^3\\
 &-480w^2-720w+120.
\end{align*}
The sixth coefficient is
\begin{equation}\label{eq:A6}
 A_6(w)=\frac{w^7}{720(1+w)^{11}}\,Q_6(w),
\end{equation}
with
\begin{align*}
 Q_6(w)={}&120w^{10}-444w^9-2920w^8+424w^7
 +16577w^6+17892w^5\\
 &-13530w^4-26400w^3-1800w^2+7200w-720.
\end{align*}
They are generated without solving any new implicit equations: one repeatedly
applies \(\mathcal L_m\) according to \cref{eq:Akoperator}.

A minimal symbolic implementation is
\begin{lstlisting}[language=Python,caption={Generating \(A_k(w)\) symbolically.}]
from sympy import symbols, diff, factorial, factor
w = symbols('w')

def L_operator(m, expr):
    return factor(w/(1+w) * diff(expr, w) - m*expr)

def A(k):
    expr = w**(k+1)
    for m in range(k):
        expr = L_operator(m, expr)
    return factor((-1)**(k+1) * expr / factorial(k+1))
\end{lstlisting}

\section{Logarithmic coefficient polynomials through order eight}\label{app:Ppolys}

Using \cref{eq:Pndef}, the next polynomials after those in
\cref{eq:Pfirst} are
\begin{align}
P_6(\ell)={}&\frac{\ell}{120}
\bigl(20\ell^5-274\ell^4+1125\ell^3-1700\ell^2
      +900\ell-120\bigr),\label{eq:P6}\\
P_7(\ell)={}&\frac{\ell}{840}
\bigl(120\ell^6-2058\ell^5+11368\ell^4-25725\ell^3
      +24500\ell^2-8820\ell+840\bigr),\label{eq:P7}\\
P_8(\ell)={}&\frac{\ell}{2520}
\bigl(315\ell^7-6534\ell^6+45962\ell^5-142149\ell^4\notag\\
&\hspace{7em}{}+205800\ell^3-135240\ell^2+35280\ell-2520\bigr).
\label{eq:P8}
\end{align}
Together with \cref{eq:Pfirst}, these determine the outer sector through
\(L^{-8}\).  Substitution into the rational \(A_k\) generates any desired
number of internal logarithmic terms in the exponentially small sectors.

\section{Independent derivation from the implicit difference equation}

For completeness, one may compute the same coefficients without Lagrange
inversion.  Substitute
\[
  d=\sum_{k=1}^{N}b_k(w)\xi^k+\bigO(\xi^{N+1})
\]
into
\[
  \left(1-\frac{d}{w}\right)(\e^{-d}+\xi)=1,
\]
expand the exponential, and equate powers of \(\xi\).  The coefficient of
\(\xi^k\) is linear in \(b_k\), because the derivative of the left-hand side
with respect to \(d\) at the origin is \(-(1+1/w)\).  Thus the recursion is
triangular.  The first three results are
\begin{align*}
 b_1(w)&=\frac{w}{1+w},\\
 b_2(w)&=\frac{w(w^2-2)}{2(1+w)^3},\\
 b_3(w)&=\frac{w(2w-1)(w^3-6w-6)}{6(1+w)^5}.
\end{align*}
Multiplying by \(w^k\) recovers \(A_k\) through
\cref{eq:bkAk}.  This route is especially convenient for proving convergence
in \(\eta\); the operator route is more convenient for symbolic generation.

\section{Reference identities for verification}

The following identities form a compact test suite for an implementation.

\begin{align*}
& W(z)\e^{W(z)}=z,\\
& \delta W(z)=\frac{W(z)}{1+W(z)},\\
& \rLambert_r(z)\e^{\rLambert_r(z)}
  +r\rLambert_r(z)=z,\\
& f_r^{-1}(z)=z-r\rLambert_r(z),\\
& \left(1-\frac{W(z)-\rLambert_r(z)}{W(z)}\right)
  \left(\exp\!\bigl(\rLambert_r(z)-W(z)\bigr)
        +r\frac{W(z)}{z}\right)=1,\\
& F\bigl(F^{-1}(z)\bigr)-z=0.
\end{align*}
For a truncation through \(z^{-N}\), the final identity should reduce to the
target remainder \(\bigO((W(z)/z)^{N+1})\), either symbolically or after
asymptotic normalization.

\clearpage
\begin{thebibliography}{99}

\bibitem{Corless1996}
R.~M. Corless, G.~H. Gonnet, D.~E.~G. Hare, D.~J. Jeffrey, and D.~E. Knuth,
``On the Lambert \(W\) function,''
\emph{Advances in Computational Mathematics} \textbf{5} (1996), 329--359.
\url{https://doi.org/10.1007/BF02124750}

\bibitem{MezoBaricz2017}
I.~Mez\H{o} and A.~Baricz,
``On the generalization of the Lambert \(W\) function,''
\emph{Transactions of the American Mathematical Society}
\textbf{369}, no.~11 (2017), 7917--7934.
\url{https://doi.org/10.1090/tran/6911}

\bibitem{DLMF}
NIST Digital Library of Mathematical Functions,
Section~4.13, ``Lambert \(W\)-Function.''
\url{https://dlmf.nist.gov/4.13}

\bibitem{FlajoletSedgewick}
P.~Flajolet and R.~Sedgewick,
\emph{Analytic Combinatorics}, Cambridge University Press, 2009.
See the discussion of Lagrange inversion and coefficient extraction.

\bibitem{Henrici}
P.~Henrici,
\emph{Applied and Computational Complex Analysis}, Vol.~1,
Wiley, 1974.

\bibitem{deBruijn}
N.~G. de Bruijn,
\emph{Asymptotic Methods in Analysis}, Dover, 1981.

\bibitem{vanDerHoeven}
J.~van der Hoeven,
\emph{Transseries and Real Differential Algebra},
Lecture Notes in Mathematics 1888, Springer, 2006.

\bibitem{ProveItDrafts}
V.~Reshetnikov,
``Series and transseries,'' in the ProveIt Fabius-function research drafts,
GitHub repository.
\url{https://github.com/VladimirReshetnikov/ProveIt/tree/main/Analysis/FabiusFunction/docs/semi-formalized-research-frontiers/drafts/series-and-transseries}
Accessed 2 September 2026.

\end{thebibliography}

\end{document}
